\documentclass[12pt,a4paper]{article}
\usepackage[utf8]{inputenc}
\usepackage[T1]{fontenc}
\usepackage[lithuanian]{babel}
\usepackage{lmodern}
\usepackage{amsmath, amssymb}
\usepackage{geometry}
\geometry{left=2cm, right=2cm, top=2.5cm, bottom=2.5cm}

\begin{document}

\begin{center}
    {\Large \textbf{Kontrolinis darbas: Daugianariai}}\\[0.3cm]
    {\large \textbf{Archimedo variantas -- Atsakymai ir sprendimai}}
\end{center}
\vspace{0.5cm}

\begin{enumerate}
    \item \textbf{Reiškinių tipai:}
    \begin{itemize}
        \item[a)] $\displaystyle \frac{5-3x}{2}$ -- \textbf{R, S} (Racionalusis sveikas). Vardiklyje nėra kintamojo, dalijama tik iš skaičiaus.
        \item[b)] $\displaystyle \frac{7}{x^3 + 4}$ -- \textbf{R, T} (Racionalusis trupmeninis). Kintamasis vardiklyje.
        \item[c)] $\displaystyle 4x^3 - \sqrt{5} x^2$ -- \textbf{R, S} (Racionalusis sveikas). Šaknis traukiama tik iš skaičiaus $5$.
        \item[d)] $\displaystyle \sqrt{x^4 - 3}$ -- \textbf{I} (Iracionalusis). Kintamasis yra po šaknies ženklu.
        \item[e)] $\displaystyle x^2 + 5^{-2}$ -- \textbf{R, S} (Racionalusis sveikas). Skaičius $5^{-2} = \frac{1}{25}$, kintamasis nėra vardiklyje.
        \item[f)] $\displaystyle \sqrt{10}x^2 + 3x$ -- \textbf{R, S} (Racionalusis sveikas). Šaknis traukiama tik iš skaičiaus $10$.
    \end{itemize}

    \vspace{0.3cm}
    \item \textbf{Vienanarių daugyba:} \\
    Pritaikome laipsnių savybes:
    $$ \left( -3x^4y^{-2} \right)^3 \cdot \left( 3x^{-5}y^4 \right)^{-2} = (-3)^3 \cdot x^{12} \cdot y^{-6} \cdot 3^{-2} \cdot x^{10} \cdot y^{-8} $$
    $$ = -27 \cdot x^{12} \cdot y^{-6} \cdot \frac{1}{9} \cdot x^{10} \cdot y^{-8} = -3 \cdot x^{12+10} \cdot y^{-6-8} = \mathbf{-3x^{22}y^{-14}} \text{ arba } \mathbf{-\frac{3x^{22}}{y^{14}}} $$
    \textbf{Paaiškinimas:} Gautas reiškinys \textbf{nėra vienanaris}, nes kintamojo $y$ laipsnio rodiklis yra neigiamas (tai reiškia dalybą iš kintamojo, o vienanariuose kintamųjų dalyti negalima).

    \vspace{0.3cm}
    \item \textbf{Skaičiavimas be skaičiuotuvo:} \\
    \begin{align*}
        \frac{156^3 + 144\left(132^2 + 2 \cdot 12 \cdot 132 + 12^2\right)}
        {156(10\sqrt{3})^2\left(12 + \frac{144^2}{156}\right)}
        &= \frac{156^3 + 144 \cdot 144^2}
        {300 \cdot 156\left(12 + \frac{144^2}{156}\right)} \\
        &= \frac{156^3 + 144^3}
        {300\left(156 \cdot 12 + 144^2\right)} \\
        &= \frac{(156+144)\left(156^2-156 \cdot 144+144^2\right)}
        {300\left(156 \cdot 12 + 144^2\right)} \\
        &= \frac{(156+144)\left(156(156-144)+144^2\right)}
        {300\left(156 \cdot 12 + 144^2\right)} \\
        &= \frac{300\left(156 \cdot 12 + 144^2\right)}
        {300\left(156 \cdot 12 + 144^2\right)} = 1.
    \end{align*}

    \vspace{0.3cm}
    \item \textbf{Formulių įrodymai:}
    \begin{itemize}
        \item[a)] $\displaystyle (x^a)^b = \underbrace{x^a \cdot x^a \cdot \dots \cdot x^a}_{b \text{ kartų}} = x^{\overbrace{a+a+\dots+a}^{b \text{ kartų}}} = \mathbf{x^{ab}}.$
        \item[b)] Atskliaudžiame dešinę pusę: \\
        $(x+y)(x^2 - xy + y^2) = x^3 - x^2y + xy^2 + x^2y - xy^2 + y^3 = \mathbf{x^3 + y^3}.$
    \end{itemize}

    \vspace{0.3cm}
    \item \textbf{Didžiausia ir mažiausia reikšmės:} \\
    Išskiriame pilnąjį kvadratą:
    $$ -2x^2 + 3x - 2 = -2\left(x^2 - \frac{3}{2}x\right) - 2 = -2\left(x^2 - 2 \cdot x \cdot \frac{3}{4} + \frac{9}{16} - \frac{9}{16}\right) - 2 $$
    $$ = -2\left( \left(x - \frac{3}{4}\right)^2 - \frac{9}{16} \right) - 2 = -2\left(x - \frac{3}{4}\right)^2 + \frac{18}{16} - 2 = -2\left(x - \frac{3}{4}\right)^2 + \frac{9}{8} - \frac{16}{8} = \mathbf{-2\left(x - \frac{3}{4}\right)^2 - \frac{7}{8}} $$
    Kadangi parabolės šakos nukreiptos žemyn, \textbf{didžiausia reikšmė yra $\mathbf{-\frac{7}{8}}$}, kai $x = \frac{3}{4}$. \\
    \textbf{Mažiausios reikšmės nėra} (artėja į $-\infty$).

    \vspace{0.3cm}
    \item \textbf{Daugianarių daugyba ir laipsnis:} \\
    Atskliaudžiame ir suprastiname reiškinį $(3x-1)^3 - (9x-1)(3x^2 - 3x + 1) + 9x\left(x - \frac{2}{3}\right)$: \\
    $ (27x^3 - 27x^2 + 9x - 1) - (27x^3 - 27x^2 + 9x - 3x^2 + 3x - 1) + 9x^2 - 6x $ \\
    $ = 27x^3 - 27x^2 + 9x - 1 - (27x^3 - 30x^2 + 12x - 1) + 9x^2 - 6x $ \\
    $ = 27x^3 - 27x^2 + 9x - 1 - 27x^3 + 30x^2 - 12x + 1 + 9x^2 - 6x $ \\
    Sutraukiame panašiuosius narius: \\
    $ = (27-27)x^3 + (-27+30+9)x^2 + (9-12-6)x + (-1+1) = \mathbf{12x^2 - 9x} $ \\
    Padauginę šį reiškinį iš $\frac{1}{x}$, gauname:
    $$ (12x^2 - 9x) \cdot \frac{1}{x} = \mathbf{12x - 9} $$
    Tai yra daugianaris, jo \textbf{laipsnis lygus 1}.

    \vspace{0.3cm}
    \item \textbf{Nežinomų koeficientų radimas:} \\
    Atskliaudžiame kairiąją lygybės pusę:
    $$ (3x-2)(2x^2+ax+1) + b = 6x^3 + 3ax^2 + 3x - 4x^2 - 2ax - 2 + b $$
    $$ = 6x^3 + (3a-4)x^2 + (3-2a)x + (b-2) $$
    Sulyginame su dešine puse ($6x^3 + 5x^2 - 5x + 3$):
    \begin{itemize}
        \item Prie $x^2$: $3a - 4 = 5 \implies 3a = 9 \implies \mathbf{a = 3}$.
        \item Laisvasis narys: $b - 2 = 3 \implies \mathbf{b = 5}$.
    \end{itemize}
    Skaičiuojame prašomo reiškinio reikšmę:
    $$ \left(-\frac{1}{3} \cdot 3\right)^{888} + \left(-\frac{1}{5} \cdot 5\right)^{777} = (-1)^{888} + (-1)^{777} = 1 + (-1) = \mathbf{0} $$

\end{enumerate}

\end{document}